%%% =======================================================================
%%% ISLS Extensions Phase 10–12 v1.0.0
%%% Phase 10: Multi-Language Forge + Full-Stack Generation
%%% Phase 11: The Alchemist — Infogenetic Signature Recombination
%%% Phase 12: Seraphic Self-Calibration
%%% Author: Sebastian Klemm
%%% Date: 18 March 2026
%%% Status: Implementation Specification for Claude Code
%%% =======================================================================
\documentclass[11pt,a4paper,twoside]{article}

\usepackage[utf8]{inputenc}
\usepackage{amsmath,amssymb,amsthm}
\usepackage{algorithm,algorithmic}
\usepackage{booktabs,longtable,tabularx}
\usepackage{enumitem}
\usepackage{geometry}
\usepackage{hyperref}
\usepackage{cleveref}
\usepackage{xcolor}
\usepackage{fancyhdr}
\usepackage{listings}
\usepackage{titlesec}

\geometry{margin=2.5cm,top=3cm,bottom=3cm,headheight=14pt}

\theoremstyle{definition}
\newtheorem{definition}{Definition}[section]
\newtheorem{axiom}{Axiom}[section]
\newtheorem{requirement}{Requirement}[section]
\newtheorem{invariant}{Invariant}[section]

\theoremstyle{plain}
\newtheorem{theorem}{Theorem}[section]

\theoremstyle{remark}
\newtheorem{remark}{Remark}[section]

\newcommand{\ISLS}{\textsf{ISLS}}

\definecolor{sec-color}{HTML}{1e3a5f}
\definecolor{dim}{HTML}{64748b}

\titleformat{\section}
  {\Large\bfseries\color{sec-color}}
  {\thesection}{1em}{}

\lstset{
  basicstyle=\ttfamily\small,
  breaklines=true,
  frame=single,
  backgroundcolor=\color{gray!5},
  rulecolor=\color{gray!30},
}

\pagestyle{fancy}
\fancyhf{}
\fancyhead[LE,RO]{\ISLS{} Extensions Phase 10--12 v1.0.0}
\fancyhead[RE,LO]{Sebastian Klemm}
\fancyfoot[C]{\thepage}

\hypersetup{
  colorlinks=true,
  linkcolor=blue!60!black,
  pdftitle={ISLS Phase 10-12},
  pdfauthor={Sebastian Klemm},
}

\begin{document}

\begin{titlepage}
\centering
\vspace*{2cm}
{\Huge\bfseries \ISLS{} Extensions\\[0.3cm]
Phase 10 -- 12\\[0.3cm]
v1.0.0}\\[1.5cm]
{\Large Phase 10: Multi-Language Forge\\
Phase 11: The Alchemist\\
Phase 12: Seraphic Self-Calibration}\\[1cm]
{\large From Rust-only to Full-Stack.\\
From Templates to Combinatorial Discovery.\\
From Manual Tuning to Self-Optimization.}\\[2cm]
{\large Sebastian Klemm}\\[0.5cm]
{\large 18 March 2026}\\[1.5cm]

\begin{abstract}
\noindent
Three phases that complete the \ISLS{} Forge as a universal software
fabrication platform.

\textbf{Phase 10} extends the Forge, Foundry, and Template system from
Rust-only to multi-language output: TypeScript/React for frontends,
Python for scripts and ML, SQL for database schemas, HTML/CSS for
static sites, YAML/Docker/CI for operations. New matrices, new
toolchain executors, new full-stack templates that produce complete
deployable applications.

\textbf{Phase 11} introduces the Alchemist: an infogenetic signature
system where every pattern, template, and crystal carries a rich
multi-dimensional signature that can be decomposed, recombined, and
mutated like genetic material. A combinatorial exploration engine (the
``slot machine'') discovers novel architectures by mixing signature
components from different domains. Keyspace shaving identifies the most
promising regions. Singularity detection flags configurations that
score anomalously high --- unexplored gold.

\textbf{Phase 12} absorbs the Seraphic Calibration Module: the Forge
tunes itself. Every forge run produces benchmark data. The Seraphic
feedback loop uses this data to adjust PMHD strategies, Oracle
prompts, template selection, and quality thresholds --- converging
toward optimal fabrication performance through a contractive
double-kick operator with Proof-of-Resonance gating.

\medskip\noindent\textbf{Status:} Implementation Specification.\\
\textbf{New Crates:} C28 (isls-multilang), C29 (isls-alchemist),
  C30 (isls-calibrate).\\
\textbf{Depends on:} C23--C27 (Forge stack), C25 (Oracle).
\end{abstract}
\end{titlepage}

\tableofcontents
\newpage


%%% =====================================================================
%%% PHASE 10
%%% =====================================================================
\part{Phase 10: Multi-Language Forge}\label{part:multilang}

\section{Purpose}\label{sec:ml:purpose}

The Forge currently produces Rust code through the RustModuleMatrix.
Real software projects are multi-language: a Rust backend needs a
TypeScript frontend, SQL migrations, Docker deployment, CI pipelines.
Phase 10 adds language-specific matrices, toolchain executors, and
full-stack templates so that a single DecisionSpec produces a
complete, multi-language, deployable project.


\section{C28: \texttt{isls-multilang}}\label{sec:ml:crate}

\begin{definition}[Language Support]\label{def:langs}
Phase 10 adds support for six additional output languages:

\begin{tabularx}{\textwidth}{l l X}
\toprule
\textbf{Language} & \textbf{Matrix} & \textbf{What It Produces} \\
\midrule
TypeScript/React & \texttt{ReactMatrix} & Components, hooks, pages,
  stores, types. JSX/TSX with Tailwind CSS. \\
Python & \texttt{PythonMatrix} & Modules, classes, FastAPI endpoints,
  CLI scripts, dataclasses, pytest tests. \\
SQL & \texttt{SqlMatrix} & DDL schemas, migrations, indexes, views,
  stored procedures, seed data. \\
HTML/CSS & \texttt{StaticMatrix} & Static pages, templates (Jinja2 or
  Handlebars), CSS with custom properties. \\
YAML/Docker & \texttt{OpsMatrix} & Dockerfiles, docker-compose,
  GitHub Actions, Kubernetes manifests, .env templates. \\
Markdown & \texttt{DocsMatrix} & README, API docs, architecture
  docs, changelogs, contribution guides. \\
\bottomrule
\end{tabularx}
\end{definition}

\begin{definition}[Toolchain Executors]\label{def:toolchains}
The Foundry (C27) gains new toolchain executors:
\begin{lstlisting}
pub enum Toolchain {
    Rust,           // cargo check, cargo test, cargo clippy
    TypeScript,     // tsc --noEmit, npm test (or vitest)
    Python,         // python -m py_compile, pytest
    Sql,            // sqlx-cli migrate (or dry parse validation)
    Docker,         // docker build --check (syntax)
    Yaml,           // YAML lint (built-in serde_yaml parse)
    Markdown,       // exists check only (always valid)
}

impl ToolchainExecutor {
    pub fn check(&self, dir: &Path, lang: &Toolchain) -> ToolchainResult;
    pub fn test(&self, dir: &Path, lang: &Toolchain) -> ToolchainResult;
    pub fn lint(&self, dir: &Path, lang: &Toolchain) -> ToolchainResult;
    pub fn fmt(&self, dir: &Path, lang: &Toolchain) -> ToolchainResult;
}
\end{lstlisting}

Toolchains are optional. If \texttt{tsc} is not installed, TypeScript
atoms are validated by the built-in syntax check (serde parse of the
AST structure). If \texttt{python} is not installed, Python atoms are
validated structurally. The system never \emph{requires} external
tools --- it degrades gracefully.
\end{definition}


\section{Full-Stack Templates}\label{sec:ml:templates}

Six new templates are added to the catalog (C26):

\begin{tabularx}{\textwidth}{l l r r X}
\toprule
\textbf{ID} & \textbf{Name} & \textbf{Atoms} & \textbf{Mol.} &
  \textbf{Languages} \\
\midrule
T11 & SaaS Starter & 22 & 6 & Rust + React + SQL + Docker + CI \\
T12 & Dashboard App & 18 & 5 & Rust + React + SQL + Docker \\
T13 & API + Docs & 12 & 3 & Rust + Markdown + Docker \\
T14 & Python ML Pipeline & 10 & 3 & Python + SQL + Docker \\
T15 & Static Site & 8 & 2 & HTML/CSS + Markdown + CI \\
T16 & Monorepo Scaffold & 16 & 5 & Rust + TS + SQL + Docker + CI \\
\bottomrule
\end{tabularx}

\subsection{T11: SaaS Starter (Full-Stack)}\label{sec:t11}

\begin{lstlisting}[title={T11 Composition Tree}]
System: saas-starter
  |
  +-- Molecule: backend (Rust)
  |     +-- Atom: api_router      [Rust, Oracle]
  |     +-- Atom: auth_service    [Rust, Oracle]
  |     +-- Atom: domain_models   [Rust, Oracle]
  |     +-- Atom: db_layer        [Rust, Oracle]
  |     +-- Atom: config          [Rust, Static]
  |     +-- Atom: main            [Rust, Static]
  |
  +-- Molecule: frontend (TypeScript/React)
  |     +-- Atom: app_shell       [TSX, Static]    -- layout, router, providers
  |     +-- Atom: auth_pages      [TSX, Oracle]    -- login, register, profile
  |     +-- Atom: dashboard_page  [TSX, Oracle]    -- main feature page
  |     +-- Atom: api_client      [TS, Derive]     -- generated from backend DTOs
  |     +-- Atom: components      [TSX, Oracle]    -- shared UI components
  |     +-- Atom: hooks           [TS, Oracle]     -- React hooks, state management
  |
  +-- Molecule: database (SQL)
  |     +-- Atom: schema          [SQL, Oracle]    -- CREATE TABLE statements
  |     +-- Atom: migrations      [SQL, Oracle]    -- migration files
  |     +-- Atom: seeds           [SQL, Oracle]    -- development seed data
  |
  +-- Molecule: ops (Docker/CI)
  |     +-- Atom: dockerfile      [Docker, Static]
  |     +-- Atom: compose         [YAML, Static]
  |     +-- Atom: ci_pipeline     [YAML, Static]  -- GitHub Actions
  |
  +-- Molecule: docs (Markdown)
        +-- Atom: readme          [MD, Derive]     -- from spec
        +-- Atom: api_docs        [MD, Derive]     -- from backend routes
\end{lstlisting}

\textbf{Interfaces:}
\begin{itemize}[nosep]
  \item Frontend \texttt{api\_client} derives from backend DTOs
        (type-safe client generation).
  \item Backend \texttt{db\_layer} consumes SQL \texttt{schema}.
  \item Docker \texttt{compose} references backend, frontend, and
        database services.
  \item CI \texttt{pipeline} runs backend tests + frontend build +
        DB migrations.
\end{itemize}

\textbf{Output:} A monorepo directory with \texttt{backend/},
\texttt{frontend/}, \texttt{database/}, \texttt{docker/},
\texttt{docs/}, \texttt{.github/workflows/}.


\section{Cross-Language Derive}\label{sec:ml:derive}

\begin{definition}[Cross-Language Derive]\label{def:crossderive}
A \emph{Derive} fill strategy can cross language boundaries:
\begin{lstlisting}
pub struct CrossDerive {
    pub source_atom: String,      // e.g., "domain_models" (Rust)
    pub target_language: String,  // e.g., "typescript"
    pub transform: DeriveTransform,
}

pub enum DeriveTransform {
    /// Extract public types and generate equivalent types
    TypeMirror,
    /// Extract API routes and generate a client
    ApiClient,
    /// Extract DB schema and generate ORM models
    OrmModels,
    /// Extract types and generate JSON Schema
    JsonSchema,
    /// Extract spec and generate documentation
    DocGen,
}
\end{lstlisting}

Example: \texttt{domain\_models} (Rust structs with serde derives) is
transformed into TypeScript interfaces automatically. No Oracle call
needed --- pure structural transformation.
\end{definition}


\section{Acceptance Tests (Phase 10)}\label{sec:ml:tests}

\begin{longtable}{l l p{5.5cm}}
\toprule
\textbf{ID} & \textbf{Name} & \textbf{Procedure} \\
\midrule
AT-ML1 & TypeScript output & Forge with ReactMatrix; verify .tsx
  files produced. \\
AT-ML2 & Python output & Forge with PythonMatrix; verify .py files. \\
AT-ML3 & SQL output & Forge with SqlMatrix; verify .sql files. \\
AT-ML4 & Cross-derive & Forge Rust backend; derive TS client; verify
  types match. \\
AT-ML5 & Full-stack T11 & Forge SaaS Starter; verify backend/,
  frontend/, database/, docker/, docs/ dirs. \\
AT-ML6 & Multi-toolchain & Run Foundry on T11 output; verify Rust
  compiles and TS type-checks (or degrades gracefully). \\
AT-ML7 & YAML validation & Forge docker-compose; verify valid YAML. \\
AT-ML8 & Monorepo scaffold & Forge T16; verify workspace Cargo.toml
  + package.json + docker-compose. \\
\end{longtable}


%%% =====================================================================
%%% PHASE 11
%%% =====================================================================
\part{Phase 11: The Alchemist}\label{part:alchemist}

\section{Purpose}\label{sec:al:purpose}

Templates are predefined. Patterns are learned from past runs. But
neither discovers \emph{novel} architectures. The Alchemist is the
combinatorial exploration engine that treats architectural decisions
as a searchable space --- mixing, recombining, and mutating structural
signatures to find configurations that no one designed on purpose but
that score anomalously high on quality metrics.

This is the ``slot machine'' for structural singularities.


\section{C29: \texttt{isls-alchemist}}\label{sec:al:crate}

\subsection{Infogenetic Signatures}\label{sec:al:sig}

\begin{definition}[Infogenetic Signature]\label{def:infosig}
Every template, pattern, and crystal carries a rich signature:
\begin{lstlisting}
pub struct InfogeneticSignature {
    /// 5D embedding (existing)
    pub embedding: FiveDState,

    /// Structural genome: ordered list of component archetypes
    pub genome: Vec<Gene>,

    /// Interaction matrix: pairwise coupling strengths between genes
    pub interactions: Vec<Vec<f64>>,

    /// Spectral fingerprint: eigenvalues of the interaction matrix
    pub spectrum: Vec<f64>,

    /// Domain tags with weights
    pub phenotype: BTreeMap<String, f64>,

    /// Quality triplet (psi, rho, omega) if measured
    pub triplet: Option<(f64, f64, f64)>,
}

pub struct Gene {
    pub id: String,
    pub kind: GeneKind,
    pub properties: BTreeMap<String, f64>,
}

pub enum GeneKind {
    Router, Handler, Service, Model, Database, Auth,
    Cache, Queue, WebSocket, Scheduler, Logger, Config,
    Frontend, Component, Hook, Store, Page,
    Migration, Schema, Seed,
    Docker, CI, Deploy,
    Custom(String),
}
\end{lstlisting}
\end{definition}

\begin{remark}
The infogenetic signature is the ``DNA'' of a software architecture.
Two architectures with similar genomes and similar interaction
matrices will have similar spectral fingerprints --- even if they were
built for different domains. This is what enables cross-domain
transfer: a pattern that works for an e-commerce backend can be
adapted for a healthcare API if their infogenetic signatures are
close.
\end{remark}

\subsection{Recombination Engine}\label{sec:al:recomb}

\begin{definition}[Recombination Operations]\label{def:recomb}
\begin{lstlisting}
pub enum AlchemyOperation {
    /// Cross two signatures: take genes from parent A and parent B
    Crossover {
        parent_a: Hash256,
        parent_b: Hash256,
        crossover_point: usize,
    },

    /// Mutate a single gene (change kind, adjust properties)
    Mutate {
        source: Hash256,
        gene_index: usize,
        mutation: GeneMutation,
    },

    /// Insert a gene from another signature into a host
    Splice {
        host: Hash256,
        donor: Hash256,
        donor_gene: usize,
        insert_at: usize,
    },

    /// Remove a gene and repair interfaces
    Excise {
        source: Hash256,
        gene_index: usize,
    },

    /// Blend: weighted average of two signatures' interaction matrices
    Blend {
        sig_a: Hash256,
        sig_b: Hash256,
        alpha: f64,  // 0.0 = pure A, 1.0 = pure B
    },
}
\end{lstlisting}
\end{definition}

\subsection{Keyspace Shaving}\label{sec:al:keyspace}

\begin{definition}[Keyspace]\label{def:keyspace}
The keyspace $K$ is the set of all possible infogenetic signatures
reachable from the current pattern memory through recombination
operations. Keyspace shaving is the process of progressively
narrowing $K$ to the most promising region:

\begin{enumerate}[nosep]
  \item Start with $K_0$ = all signatures in pattern memory.
  \item Select the operation with highest expected quality gain:
        $q(g) = \frac{H(K) - H(K \cap g)}{\mathrm{risk}(g) + \varepsilon}$
  \item Apply the operation: $K_{t+1} = K_t \cap g^*$.
  \item Repeat until $\mathrm{Var}(K_t) \le \tau$ or budget exhausted.
\end{enumerate}
\end{definition}

\subsection{Singularity Detection}\label{sec:al:singularity}

\begin{definition}[Singularity]\label{def:singularity}
A singularity is a discovered signature that scores $> 2\sigma$ above
the mean quality of all patterns in its spectral neighborhood:
\[
  \mathrm{singularity}(\mathbf{s}) =
  \frac{Q(\mathbf{s}) - \mu_{\mathrm{neighbors}}}
       {\sigma_{\mathrm{neighbors}}} > 2
\]
where $Q$ is the quality functional and the neighborhood is defined
by spectral fingerprint distance $< \theta_{\mathrm{spec}}$.

A singularity represents unexplored gold: an architecture that is
significantly better than anything the system has seen in that region
of signature space.
\end{definition}

\subsection{The Alchemy Lab (Studio Integration)}\label{sec:al:studio}

\begin{lstlisting}[title={New Studio View: The Alchemy Lab}]
+--------------------------------------------------+
|  ALCHEMY LAB                                       |
+--------------------------------------------------+
|                                                    |
|  +-- Signature Viewer ----------------------------+ |
|  | Selected: pattern p-a7b3 (REST API)            | |
|  | Genome: [Router, Auth, Service, Model, DB]     | |
|  | Spectrum: [0.91, 0.73, 0.45, 0.22, 0.08]      | |
|  | Quality: psi=0.82 rho=0.76 omega=0.91          | |
|  +------------------------------------------------+ |
|                                                    |
|  +-- Recombination Panel -------------------------+ |
|  | Parent A: [p-a7b3 (REST API)     v]            | |
|  | Parent B: [p-f2d1 (WebSocket)    v]            | |
|  | Operation: [Crossover v] at gene 3             | |
|  |                                                 | |
|  | [ RECOMBINE ] -> Preview child signature        | |
|  | [ FORGE CHILD ] -> Materialize as real project  | |
|  +------------------------------------------------+ |
|                                                    |
|  +-- Keyspace Map (canvas) -----------------------+ |
|  | [2D projection of signature space]              | |
|  | Dots = patterns, Stars = singularities          | |
|  | Color = quality score                           | |
|  +------------------------------------------------+ |
|                                                    |
|  +-- Singularity Feed ----------------------------+ |
|  | ! SINGULARITY: REST+WebSocket+Queue hybrid      | |
|  |   Score: 0.94 (2.3 sigma above neighborhood)   | |
|  |   [Inspect] [Forge] [Save as Template]          | |
|  +------------------------------------------------+ |
+--------------------------------------------------+
\end{lstlisting}


\section{Acceptance Tests (Phase 11)}\label{sec:al:tests}

\begin{longtable}{l l p{5.5cm}}
\toprule
\textbf{ID} & \textbf{Name} & \textbf{Procedure} \\
\midrule
AT-AL1 & Signature creation & Create InfogeneticSignature from a
  template; verify genome and spectrum. \\
AT-AL2 & Crossover & Cross two signatures; verify child has genes
  from both parents. \\
AT-AL3 & Mutation & Mutate a gene; verify signature changes. \\
AT-AL4 & Splice & Insert a gene from donor; verify interfaces
  repaired. \\
AT-AL5 & Keyspace shaving & Shave from 100 signatures; verify
  variance decreases. \\
AT-AL6 & Singularity detection & Inject a high-scoring outlier;
  verify detected as singularity. \\
AT-AL7 & Forge from recombination & Recombine two signatures; forge
  the child; verify working project. \\
AT-AL8 & Spectral fingerprint & Verify two similar templates have
  similar spectral fingerprints. \\
\end{longtable}


%%% =====================================================================
%%% PHASE 12
%%% =====================================================================
\part{Phase 12: Seraphic Self-Calibration}\label{part:seraphic}

\section{Purpose}\label{sec:se:purpose}

The Forge has many knobs: PMHD tick count, opposition strength,
quality thresholds, Oracle temperature, template match threshold,
retry budget, fill strategy weights. Currently these are
hand-configured. Phase 12 makes the Forge tune itself.

The Seraphic Calibration Shell (absorbed from the Q$\otimes$DASH SCS
document) wraps the Forge in a feedback loop: every forge run
produces benchmark data (compilation success rate, test pass rate,
Oracle rejection rate, token usage, autonomy ratio). This data feeds
a contractive double-kick operator that adjusts the Forge's
configuration toward a fixpoint attractor --- the optimal operating
point for the current domain and workload.


\section{C30: \texttt{isls-calibrate}}\label{sec:se:crate}

\begin{definition}[Forge Configuration Space]\label{def:configspace}
The configuration space $C$ includes:
\begin{lstlisting}
pub struct ForgeCalibrationState {
    // PMHD
    pub pmhd_ticks: u64,
    pub pmhd_pool_size: usize,
    pub pmhd_opposition_strength: f64,
    pub pmhd_mutation_rate: f64,
    pub pmhd_strategy: SeedStrategy,

    // Oracle
    pub oracle_temperature: f64,
    pub oracle_max_tokens: usize,
    pub oracle_max_retries: usize,

    // Templates
    pub template_match_threshold: f64,

    // Quality
    pub quality_thresholds: QualityMetrics,

    // Foundry
    pub foundry_max_attempts: usize,
}
\end{lstlisting}
\end{definition}

\begin{definition}[Performance Triplet]\label{def:perftriplet}
For each forge run, a performance triplet is computed:
\begin{lstlisting}
pub struct ForgeTriplet {
    pub psi: f64,   // quality: test pass rate * constraint satisfaction
    pub rho: f64,   // stability: 1 - (variance across retries)
    pub omega: f64,  // efficiency: 1 / (tokens_used / baseline)
}
\end{lstlisting}
\end{definition}

\begin{definition}[Seraphic Feedback]\label{def:feedback}
After each forge run, the benchmark vector $u$ is encoded into the
calibration field:
\begin{lstlisting}
pub struct CalibrationField {
    pub state: Vec<f64>,       // M(t) field state
    pub history: Vec<ForgeTriplet>, // past triplets
}

impl CalibrationField {
    pub fn update(&mut self, triplet: &ForgeTriplet, injection: &[f64]) {
        // M(t+1) = Norm(alpha * M(t) + gamma * I_t)
        for i in 0..self.state.len() {
            self.state[i] = (0.9 * self.state[i]
                           + 0.1 * injection[i]).clamp(0.0, 1.0);
        }
        self.history.push(triplet.clone());
    }
}
\end{lstlisting}
\end{definition}

\begin{definition}[Double-Kick Operator]\label{def:doublekick}
The configuration update operator $T = \Phi_V \circ \Phi_U$:
\begin{lstlisting}
pub fn double_kick(
    config: &mut ForgeCalibrationState,
    triplet: &ForgeTriplet,
    field: &CalibrationField,
) {
    // Phi_U: quality-ascending update
    if triplet.psi < target_psi {
        config.pmhd_ticks += 100;          // more drill
        config.oracle_max_retries += 1;     // more attempts
    }
    if triplet.psi > target_psi * 1.1 {
        config.pmhd_ticks = (config.pmhd_ticks - 50).max(100);
        config.oracle_max_retries = (config.oracle_max_retries - 1).max(1);
    }

    // Phi_V: stability/efficiency stabilization
    if triplet.rho < target_rho {
        config.pmhd_opposition_strength += 0.05;
        config.oracle_temperature = (config.oracle_temperature - 0.05).max(0.0);
    }
    if triplet.omega < target_omega {
        config.oracle_max_tokens = (config.oracle_max_tokens - 256).max(1024);
        config.template_match_threshold -= 0.05; // use templates more
    }

    // Clamp all values to admissible ranges
    config.clamp();
}
\end{lstlisting}
\end{definition}

\begin{definition}[Proof-of-Resonance for Config]\label{def:porcfg}
A config update is accepted only if:
\begin{enumerate}[nosep]
  \item $\psi(c') \ge \psi(c)$ (no quality regression),
  \item $\rho(c')$ within tolerance of $\rho(c)$,
  \item $\omega(c')$ above minimum threshold,
  \item Field $M(t)$ shows positive correlation with the new config's
        injection.
\end{enumerate}
If PoR fails, the config update is rejected and the system continues
with the previous configuration.
\end{definition}

\begin{definition}[CRI Regime Change]\label{def:cri}
If the global quality $J(t) = \psi \cdot \rho \cdot \omega$ stagnates
for $N$ consecutive runs, a CRI impulse triggers:
\begin{itemize}[nosep]
  \item Switch PMHD strategy (e.g., Hybrid $\to$ Evolutionary).
  \item Increase Oracle temperature temporarily (explore more).
  \item Reset template match threshold to default.
  \item Log the regime change as a calibration event.
\end{itemize}
After the CRI impulse, the double-kick resumes in the new region.
\end{definition}


\section{Studio Integration}\label{sec:se:studio}

\begin{lstlisting}[title={Enhanced Dashboard: Calibration Panel}]
+-- Forge Calibration ----------------------------------+
| Triplet: psi=0.87 rho=0.91 omega=0.76                |
| Global J: 0.60  (trend: +0.03 over last 10 runs)     |
| Regime: Hybrid | Calibration step: 47                 |
| Last CRI: step 23 (Greedy -> Hybrid)                  |
|                                                        |
| [Triplet chart over time: psi/rho/omega lines]         |
| [Config evolution: ticks, temperature, threshold]      |
+--------------------------------------------------------+
\end{lstlisting}


\section{Acceptance Tests (Phase 12)}\label{sec:se:tests}

\begin{longtable}{l l p{5.5cm}}
\toprule
\textbf{ID} & \textbf{Name} & \textbf{Procedure} \\
\midrule
AT-SC1 & Triplet computation & Run forge; verify triplet
  $(\psi,\rho,\omega) \in [0,1]^3$. \\
AT-SC2 & Field update & Run 5 forges; verify field state changes
  monotonically toward better configs. \\
AT-SC3 & Double-kick & Start with suboptimal config; verify
  config adjusts after 3 runs. \\
AT-SC4 & PoR rejection & Mock a config that regresses quality;
  verify rejected and old config restored. \\
AT-SC5 & CRI trigger & Stagnate J for 10 runs; verify CRI
  impulse triggers strategy switch. \\
AT-SC6 & Config persistence & Calibrate; restart; verify
  calibrated config loaded from store. \\
\end{longtable}


%%% =====================================================================
\part{Dependency Graph and Handoff}\label{part:handoff}

\section{New Crates}

\begin{tabularx}{\textwidth}{l l X}
\toprule
\textbf{Crate} & \textbf{ID} & \textbf{Depends On} \\
\midrule
\texttt{isls-multilang} & C28 & C22 (ArtifactIR), C23 (Forge),
  C26 (Templates), C27 (Foundry) \\
\texttt{isls-alchemist} & C29 & C22, C23, C25 (Oracle),
  C26 (Templates) \\
\texttt{isls-calibrate} & C30 & C23 (Forge), C25 (Oracle),
  C17 (Store) \\
\bottomrule
\end{tabularx}

\section{Test Count}

\begin{tabular}{l r r}
\toprule
\textbf{Phase} & \textbf{New Tests} & \textbf{Cumulative} \\
\midrule
Phase 1--9 & 326 & 326 \\
Phase 10 (C28) & 8 & 334 \\
Phase 11 (C29) & 8 & 342 \\
Phase 12 (C30) & 6 & $\ge$ 348 \\
\bottomrule
\end{tabular}

\section{Completion Criteria}

\begin{enumerate}[nosep]
  \item Forge produces TypeScript, Python, SQL, HTML, YAML, Docker,
        Markdown output.
  \item Full-stack templates (T11--T16) produce multi-language
        projects.
  \item Cross-language derive generates TS types from Rust structs.
  \item Foundry validates with language-appropriate toolchains.
  \item Infogenetic signatures computed for all patterns and
        templates.
  \item Recombination produces viable child signatures.
  \item Singularity detection identifies outlier architectures.
  \item Alchemy Lab view in Studio with keyspace visualization.
  \item Seraphic calibration adjusts forge config across runs.
  \item PoR gates config changes; CRI triggers regime switches.
  \item Total tests $\ge 348$, zero warnings.
  \item \textbf{30 Crates. Every language. Self-tuning.
        Architecture DNA.}
\end{enumerate}


\vfill
\begin{center}
\rule{0.5\textwidth}{0.4pt}\\[0.5cm]
{\small Generated for \ISLS{} Extensions Phase 10--12 v1.0.0.\\[0.2cm]
Phase 10: The Forge speaks every language.\\
Phase 11: The Alchemist discovers what no one designed.\\
Phase 12: The system tunes itself.\\[0.3cm]
Infogenetic signatures. Combinatorial exploration. Fixpoint convergence.\\
Every pattern has DNA. Every recombination is an experiment.\\
Every singularity is unexplored gold.}
\end{center}

\end{document}
